\begin{textsample}{NEG Dim 2 – global\_south\_2025\_02 – Score 1.00 – global\_south\_2025\_02/120982219.txt}  \label{ex:f2_neg_028}
Trump ' s Gaza plan : There is more to a home than real estate Summary The US president ' s ' Riviera ' proposal for this war-ravaged Palestinian territory has met global rejection for good reason.
Now the world must speak in unison to remind America of the perils of a colonial project This is a Mint Premium article gifted to you.
Subscribe to enjoy similar stories.
Even by his own standards, US President Donald Trump ' s announcement on Gaza this week was a head-spinner.
America, he said, would take over that war-ravaged strip of land, move its 2-million-odd people out, secure it with US troops if necessary, and then build a grand"Riviera of the Middle East"by the Mediterranean Sea where anybody from around the world could eventually come and live.
His declaration was met with global consternation.
Nearly every country rejected the idea, including all other members of the UN Security Council.
Brics nations have weighed in with their criticism, as have countries Trump may @ @ @ @ @ @ @ @ @ @, like Saudi Arabia, Jordan, Egypt and Turkey.
Leaders have called the plan bizarre, absurd, inhumane and unworkable.
Gaza ' s Palestinians, evicted from their homes as Israel flattened their homeland with bombs, must be allowed to return home, leaders argued.
London said nothing should be done that unsettles the fragile Hamas-Israel ceasefire.
Ant? nio Guterres, the UN ' s secretary-general, warned against"any form of ethnic cleansing."To be sure, White House officials appeared to backpedal on the proposal a day after it was made, saying that Palestinian displacement would be temporary, although Trump had said"either"when asked if he meant an interim or permanent eviction.
The White House press secretary said Trump believes a US role in rebuilding Gaza would"ensure stability in the region,"adding,"That does not mean boots on the ground in Gaza."Whether Trump plucked the idea out of a hat or whether it has been in the works is not the @ @ @ @ @ @ @ @ @ @.
What matters is that it can not be brushed aside as just another Trump throwaway.
This is why world leaders rejected it so promptly.
An incendiary plan for a volatile region could set it up for another flashpoint, with risks to countries seen approving of it.
Second, while Trump appears to favour territorial expansion, a big shift in itself, such a brazen colonial project would deal a blow to the stability assurance of Pax Americana, which may spell perils elsewhere too.
A world order where might is right suits nobody.
Even if Trump saw the strip as a deal-making chip, with a maximal stance taken tactically in public to buy Israel ' s leader some time to survive a far-right political revolt over Gaza being returned to its inhabitants, his Riviera pitch is far too reckless.
Words do send out signals, especially in a context fraught with America ' s military record in the region.
What can the rest of the world do?
Speak in one voice, for a start, @ @ @ @ @ @ @ @ @ @ to everybody.
This is not an era of war, as Indian Prime Minister Narendra Modi has said.
Nor is it one of forced evictions.
It is one thing to disrupt trade patterns, quite another to redraw the world map -- in blood or otherwise.
We must also reiterate that the Palestinian cause can not be wished away.
The people who lived on that land before the 1948 carve-out of Israel have endured much suffering.
Whatever Trump may be bent on doing, a two-state solution remains the optimal way out of this crisis.
Upwards of 47,000 people -- mostly women and children -- may have been killed by Israel ' s Gaza war in retaliation to the Hamas terror attack of 7 October 2023.
Of course Gaza must be rebuilt, but by Palestinians themselves, aided by others, not by a colonial force.

% matched lemmas: 
\end{textsample}
